\documentclass[UTF8,11pt,a4paper]{ctexart}

% =========================
% 基础宏包
% =========================
\usepackage{amsmath,amssymb,amsfonts,amsthm,mathtools,bm}
\usepackage{geometry}
\usepackage{booktabs,multirow,array,longtable}
\usepackage{graphicx}
\usepackage{xcolor}
\usepackage{enumitem}
\usepackage{hyperref}
\usepackage[nameinlink,capitalize]{cleveref}
\usepackage[numbers,sort&compress]{natbib}

\geometry{margin=1in}
\hypersetup{
  colorlinks=true,
  linkcolor=blue!55!black,
  citecolor=blue!55!black,
  urlcolor=blue!55!black,
  pdftitle={谱梯度更新的激活扰动几何},
  pdfauthor={TBD}
}

% =========================
% 定理环境
% =========================
\newtheorem{theorem}{定理}[section]
\newtheorem{proposition}[theorem]{命题}
\newtheorem{lemma}[theorem]{引理}
\newtheorem{corollary}[theorem]{推论}
\theoremstyle{definition}
\newtheorem{definition}[theorem]{定义}
\newtheorem{assumption}[theorem]{假设}
\newtheorem{remark}[theorem]{注记}

% =========================
% 常用命令
% =========================
\newcommand{\R}{\mathbb{R}}
\newcommand{\E}{\mathbb{E}}
\newcommand{\Tr}{\operatorname{tr}}
\newcommand{\rank}{\operatorname{rank}}
\newcommand{\diag}{\operatorname{diag}}
\newcommand{\vecop}{\operatorname{vec}}
\newcommand{\op}{\mathrm{op}}
\newcommand{\srank}{\operatorname{srank}}
\newcommand{\nrank}{\operatorname{nrank}}
\newcommand{\ssrank}{\operatorname{ssrank}}
\newcommand{\SpecGrad}{\mathrm{SpecGrad}}
\newcommand{\Loss}{\mathcal{L}}
\newcommand{\D}{\mathrm{d}}
\newcommand{\TBD}{\textcolor{red!70!black}{\textbf{TBD}}}

\title{谱梯度更新的激活扰动几何：\\局部理论与更新谱塑形实验证据}
\author{TBD}
\date{\today}

\begin{document}
\maketitle

\begin{abstract}
谱梯度方法（spectral gradient methods）通过重塑矩阵更新的奇异值来改变优化动态。典型例子是极分解更新
\(G=U\Sigma V^\top\mapsto UV^\top\)，它保留梯度奇异向量，但将非零奇异值拉平。已有理论通常从逐层角度解释这类方法：当某个参数块的梯度核秩（nuclear rank）大于输入激活的稳定秩（stable rank）时，谱更新可能优于欧几里得梯度步。本文将这个观点整理为一个更一般的\emph{激活扰动几何}框架。

我们首先证明：在一阶线性化目标下，梯度形状的更新是 Frobenius 范数预算下的最优方向，而极分解/谱更新是 operator 范数预算下的最优方向。这解释了为什么高秩更新并不天然带来更大损失下降。其次，我们把参数更新对隐藏激活的影响写成 \(\Delta W A\)，并指出谱更新提供的是 per-sample/operator-norm 意义下的扰动控制，但它可能增加 batch-level activation spreading。进一步地，若该层扰动通过下游 Jacobian \(B\) 传播到模型输出，则真正相关的量不是 \(\srank(A)\)，而是
\[
\ssrank(B,A)=
\frac{\sum_i \sigma_i(B)^2\sigma_i(A)^2}{\sigma_1(B)^2\sigma_1(A)^2},
\]
这是映射 \(D\mapsto BDA\) 在谱范数与 Frobenius 范数几何下的精确敏感度比。

实验部分目前采用 E11 控制实验作为局部机制证据：Muon-like 极分解更新稳定地产生更平、更高秩的 update spectrum；one-step 损失下降与梯度--更新内积高度相关；flat/polar 谱分配在 operator-norm 预算下胜出，但在 Frobenius 预算下失败；direct activation perturbation diagnostics 显示 Muon 在 matched update size 下可降低 operator/per-sample worst-case perturbation，但不统一降低 batch-Frobenius movement；神经网络 sanity checks 进一步说明高秩更新不自动意味着更好的一阶进展。尚未进行的 downstream Jacobian 与长尾小 batch 实验在本文中保留为待完成实验。
\end{abstract}


\tableofcontents
\newpage

\section{引言}

深度网络中大量参数天然是矩阵：全连接层、MLP 投影、嵌入矩阵、卷积核的矩阵化表示、以及 Transformer 中的投影矩阵。标准一阶优化方法通常在欧几里得几何或坐标几何中更新这些矩阵。谱梯度方法则不同：它们直接作用于矩阵梯度的奇异值结构。

设某个矩阵参数块的梯度为
\[
G=U\Sigma V^\top.
\]
一般谱梯度更新可以写成
\[
D_f(G)=Uf(\Sigma)V^\top,
\]
其中 \(f\) 对奇异值进行重标定。最典型的极分解更新为
\[
D_{\mathrm{polar}}(G)=UV^\top.
\]
它保留梯度奇异向量，但将非零奇异值全部设置为同一尺度。

近期工作已经从不同角度解释了这种更新。Davis 和 Drusvyatskiy 提出，当梯度的 squared nuclear-to-Frobenius ratio 大于输入激活稳定秩时，谱更新的一步下降界可能优于欧几里得梯度步\citep{davis2026spectral}. Chen、Li 和 Liu 则将 Muon 放入 Lion-\(\mathcal K\) 框架，指出 Muon 对应核范数选择，并与谱范数约束优化相关\citep{chen2025muon}. 这些工作表明，谱更新不是单纯的数值技巧，而是对应一种不同的矩阵几何。

本文从另一个角度组织这一现象：\emph{激活扰动几何}。给定某层输入激活矩阵
\[
A_{\ell-1}=[h_{\ell-1}(x_1),\ldots,h_{\ell-1}(x_B)],
\]
该层 pre-activation 为
\[
Z_\ell=W_\ell A_{\ell-1}.
\]
若参数更新为
\[
W_\ell^+=W_\ell-D_\ell,
\]
则直接 pre-activation 扰动为
\[
\Delta Z_{\ell,\mathrm{direct}}=-D_\ell A_{\ell-1}.
\]
因此，讨论谱更新是否稳定或有效，不能只看 \(D_\ell\) 的 rank 或 stable rank，而必须看
\[
\|D_\ell A_{\ell-1}\|
\]
以及该扰动通过后续网络传播后的函数影响
\[
\|B_\ell D_\ell A_{\ell-1}\|.
\]

本文的中心观点是：
\[
\boxed{\text{谱梯度更新是 operator-norm 几何下自然的一阶方向，但其优化收益取决于激活与下游 Jacobian 的几何。}}
\]
换言之，高秩/平谱更新只是机制事实；它是否带来损失下降或激活稳定，需要由 gradient alignment、activation stable rank、downstream sensitivity 共同决定。

\section{相关工作}

\paragraph{谱梯度更新与 Muon.}
谱梯度方法将矩阵梯度的奇异值替换为某种变换后的谱。极分解更新 \(G=U\Sigma V^\top\mapsto UV^\top\) 是 operator-norm 约束下一阶最速下降方向。Muon 可以视为这种谱更新思想的实际优化器实现或近似实现。\citet{chen2025muon} 将 Muon 解释为 Lion-\(\mathcal K\) 框架中核范数对应的情形，并分析其与谱范数约束优化之间的联系。

\paragraph{层级 nuclear-rank vs stable-rank 条件.}
\citet{davis2026spectral} 提出，谱更新的优势可由逐层条件解释：对于梯度矩阵 \(G_\ell\) 和输入激活矩阵 \(A_{\ell-1}\)，当
\[
\frac{\|G_\ell\|_*^2}{\|G_\ell\|_F^2}
>
\frac{\|A_{\ell-1}\|_F^2}{\|A_{\ell-1}\|_{\op}^2}
\]
时，谱更新相对于欧几里得更新具有更好的 one-step 保证下降。本文将这个条件解释为 \(B=I\) 的特殊情况，并引入下游敏感度 \(B\) 后的 sandwiched stable rank。

\paragraph{Stable rank normalization.}
稳定秩也出现在泛化与 Lipschitz 控制中。\citet{sanyal2020stable} 提出 Stable Rank Normalization，通过控制线性层稳定秩和经验 Lipschitz 行为来改善泛化。本文的 stable rank 用法不同：它不是直接作为正则项，而是作为局部扰动几何中的敏感度比。

\paragraph{E11 实验套件定位.}
本文使用的已有实验来自 E11 condition geometry 分支。该实验套件的问题设定是：Muon 是否作为 geometry-shaping optimizer 改变更新矩阵谱几何，以及这种几何是否解释 one-step 或短程优化进展。其当前 thesis 是：Muon 会产生更平、更高秩的 update spectrum，但是否有利取决于任务、层、范数约束、horizon 和调参。该实验套件在本文中作为局部谱几何机制证据，而非完整优化器性能证据。

\section{问题设定与记号}

设模型参数由矩阵块组成：
\[
\theta=(W_1,\ldots,W_L),\qquad W_\ell\in\R^{m_\ell\times n_\ell}.
\]
对一个 batch，模型输出为
\[
F(\theta)\in\R^{c\times B},
\]
损失为
\[
\Loss(\theta)=\phi(F(\theta)).
\]
令
\[
G_\ell=\nabla_{W_\ell}\Loss(\theta)
\]
为第 \(\ell\) 个矩阵块的梯度。

若某层输入 activation matrix 为
\[
A_{\ell-1}=[h_{\ell-1}(x_1),\ldots,h_{\ell-1}(x_B)]\in\R^{n_\ell\times B},
\]
则该层 pre-activation 为
\[
Z_\ell=W_\ell A_{\ell-1}.
\]
设正下降更新为
\[
D_\ell=W_\ell-W_\ell^+.
\]
那么固定输入激活时的直接 pre-activation 变化为
\[
\Delta Z_{\ell,\mathrm{direct}}=-D_\ell A_{\ell-1}.
\]
如果激活函数 \(\sigma\) 是 \(L_\sigma\)-Lipschitz，则
\[
\|\Delta H_\ell\|_F
\le
L_\sigma\|D_\ell A_{\ell-1}\|_F.
\]
因此，本文将 \(D_\ell A_{\ell-1}\) 视为最基本的 activation perturbation object。

对矩阵 \(M\)，记
\[
\|M\|_{\op}=\sigma_1(M),\qquad
\|M\|_F=\left(\sum_i\sigma_i(M)^2\right)^{1/2},
\qquad
\|M\|_*=\sum_i\sigma_i(M).
\]
稳定秩定义为
\[
\srank(M)=\frac{\|M\|_F^2}{\|M\|_{\op}^2},
\]
梯度核秩定义为
\[
\nrank(G)=\frac{\|G\|_*^2}{\|G\|_F^2}.
\]

\section{局部一阶谱几何}

本节说明为何 polar/SpecGrad 更新自然对应 operator-norm 几何，而非 Frobenius 几何。

\subsection{SpecGrad 家族}

设
\[
G=U\Sigma V^\top
\]
为某个矩阵梯度。定义谱梯度更新族：
\[
D_f(G)=Uf(\Sigma)V^\top.
\]
特殊情形包括：
\[
f(\sigma)=\sigma \quad \Rightarrow \quad D_f(G)=G,
\]
\[
D_F(G)=G/\|G\|_F,
\]
以及 polar 方向
\[
f(\sigma)=1_{\sigma>0},\qquad D_{\mathrm{polar}}(G)=UV^\top.
\]

\subsection{Frobenius 与 operator-norm 约束下的一阶最优性}

在当前参数点，对单个矩阵块考虑局部展开：
\[
\Loss(W-D)-\Loss(W)
= -\langle G,D\rangle+O(\|D\|^2).
\]
因此，一阶目标是最大化
\[
\langle G,D\rangle.
\]

\begin{proposition}[谱分配边界]
设 \(G=U\operatorname{diag}(\sigma)V^\top\)。
\begin{enumerate}[label=(\roman*)]
  \item 在 Frobenius 球 \(\|D\|_F\le r\) 上，唯一最优方向为
  \[
  D_F^\star=rG/\|G\|_F,
  \]
  最优一阶下降为
  \[
  r\|G\|_F.
  \]
  \item 在 operator-norm 球 \(\|D\|_{\op}\le \eta\) 上，一个最优方向为
  \[
  D_{\op}^\star=\eta UV^\top,
  \]
  最优一阶下降为
  \[
  \eta\|G\|_*.
  \]
\end{enumerate}
\end{proposition}

\begin{proof}
Frobenius 情况由 Cauchy--Schwarz 给出：
\[
\langle G,D\rangle\le \|G\|_F\|D\|_F\le r\|G\|_F,
\]
等号在 \(D=rG/\|G\|_F\) 处取得。

Operator-norm 情况由 von Neumann trace inequality 给出：
\[
\langle G,D\rangle
\le \sum_i\sigma_i(G)\sigma_i(D)
\le \eta\sum_i\sigma_i(G)
=\eta\|G\|_*.
\]
取 \(D=\eta UV^\top\) 达到等号。
\end{proof}

该命题解释了一个关键事实：
\[
\boxed{\text{polar/flat spectrum 不是普适最优，而是 operator-norm 预算下的一阶最优。}}
\]
如果公平比较的更新预算是 Frobenius 范数，那么梯度形状 \(G/\|G\|_F\) 才是一阶最优；如果公平比较的是 operator norm，那么 polar 更新才是一阶最优。

\section{激活扰动几何}

本节把上一节的一阶更新几何转换为 activation perturbation 的语言。

考虑单层线性映射
\[
Z=WA,
\]
其中 \(A\) 是输入激活矩阵。更新 \(W^+=W-\eta D\) 后，直接 pre-activation 扰动为
\[
\Delta Z=-\eta DA.
\]

\subsection{Polar 更新的 operator-norm 控制}

若
\[
D=D_{\mathrm{polar}}=UV^\top,
\]
则
\[
\|D\|_{\op}=1.
\]
因此对任意单样本 activation vector \(a\)，有
\[
\|Da\|_2\le \|a\|_2.
\]
若步长为 \(\eta\)，则
\[
\|\Delta z(x)\|_2\le \eta\|h_{\ell-1}(x)\|_2.
\]
这说明 polar/SpecGrad 更新给出了 per-sample worst-case activation perturbation bound。

对整个 batch，有
\[
\|DA\|_F\le \|D\|_{\op}\|A\|_F=\|A\|_F.
\]
相反，Frobenius-normalized gradient direction
\[
D_F=G/\|G\|_F
\]
满足
\[
\|D_FA\|_F\le \|D_F\|_F\|A\|_{\op}=\|A\|_{\op}.
\]
由于
\[
\|A\|_F=\sqrt{\srank(A)}\|A\|_{\op},
\]
高稳定秩 activation 下，polar 更新可能导致更大的 batch-level activation movement。

因此，本文区分两种稳定性：
\[
\boxed{\text{operator/per-sample activation stability}}
\]
与
\[
\boxed{\text{batch-Frobenius activation stability}.}
\]
SpecGrad 更自然地控制前者，而非后者。

\subsection{层级谱更新优势条件}

设 loss 对 \(Z=WA\) 是 \(\beta\)-smooth。则
\[
\Loss(W-\eta D)
\le
\Loss(W)-\eta\langle G,D\rangle+\frac{\beta\eta^2}{2}\|DA\|_F^2.
\]

对 Frobenius-normalized direction \(D_F=G/\|G\|_F\)，有
\[
\langle G,D_F\rangle=\|G\|_F,
\qquad
\|D_FA\|_F\le \|A\|_{\op}.
\]
因此最优步长下的保证下降为
\[
\operatorname{Dec}_F\gtrsim
\frac{\|G\|_F^2}{2\beta\|A\|_{\op}^2}.
\]

对 polar direction \(D_S=UV^\top\)，有
\[
\langle G,D_S\rangle=\|G\|_*,
\qquad
\|D_SA\|_F\le \|A\|_F.
\]
因此
\[
\operatorname{Dec}_S\gtrsim
\frac{\|G\|_*^2}{2\beta\|A\|_F^2}.
\]

于是 spectral update 更优的充分条件为
\[
\frac{\|G\|_*^2}{\|G\|_F^2}
>
\frac{\|A\|_F^2}{\|A\|_{\op}^2},
\]
也就是
\[
\boxed{\nrank(G)>\srank(A).}
\]

\section{下游敏感度与 Sandwiched Stable Rank}

单层条件只考虑 \(DA\) 对该层输出的扰动。但在深层网络中，该扰动还要经过后续层传播。在线性化下，可写成
\[
\Delta F\approx BDA,
\]
其中 \(B\) 是从该层输出到最终模型输出的 downstream Jacobian。

\subsection{精确敏感度公式}

定义
\[
K_F^2=\sup_{\|D\|_F\le 1}\|BDA\|_F^2,
\]
以及
\[
K_S^2=\sup_{\|D\|_{\op}\le 1}\|BDA\|_F^2.
\]

\begin{theorem}[Sandwiched block 的精确敏感度]
对映射 \(J(D)=BDA\)，有
\[
K_F^2=\|B\|_{\op}^2\|A\|_{\op}^2,
\]
且
\[
K_S^2=\sum_i\sigma_i(B)^2\sigma_i(A)^2,
\]
其中奇异值按降序排列，维度不足处以 \(0\) 补齐。
\end{theorem}

\begin{proof}
Frobenius 情况使用 vectorization：
\[
\vecop(BDA)=(A^\top\otimes B)\vecop(D).
\]
因此
\[
K_F=\|A^\top\otimes B\|_{\op}=\|A\|_{\op}\|B\|_{\op}.
\]

Operator-norm 情况令
\[
B=U_B\Sigma_BV_B^\top,
\qquad
A=U_A\Sigma_AV_A^\top.
\]
写
\[
Z=V_B^\top D U_A,
\]
则 \(\|Z\|_{\op}\le \|D\|_{\op}\)，且
\[
\|BDA\|_F^2
=\|\Sigma_B Z\Sigma_A\|_F^2.
\]
根据 von Neumann rearrangement inequality，
\[
\sup_{\|Z\|_{\op}\le 1}\|\Sigma_B Z\Sigma_A\|_F^2
=
\sum_i\sigma_i(B)^2\sigma_i(A)^2.
\]
取使奇异向量对齐的 partial isometry 可达等号。
\end{proof}

\subsection{Sandwiched stable rank}

定义
\[
\boxed{
\ssrank(B,A)=
\frac{\sum_i\sigma_i(B)^2\sigma_i(A)^2}
{\sigma_1(B)^2\sigma_1(A)^2}.
}
\]
则
\[
\frac{K_S^2}{K_F^2}=\ssrank(B,A).
\]
并且
\[
1\le \ssrank(B,A)
\le
\min\{\srank(B),\srank(A)\}.
\]
特别地，当 \(B=I\) 时，
\[
\ssrank(I,A)=\srank(A).
\]
因此逐层 activation stable rank 是 downstream-aware sensitivity ratio 的特殊情形。

对应的下游修正谱更新条件为
\[
\boxed{\nrank(G)>\ssrank(B,A).}
\]
该条件比 \(\nrank(G)>\srank(A)\) 更精细，因为它保留了下游 Jacobian 的奇异值谱。

\section{已有 E11 实验证据}
\label{sec:e11}

本节整理当前已有实验结果。它们支持本文的局部谱几何机制，并包含 direct activation perturbation 的初步证据；downstream Jacobian 理论仍未被直接验证。

\subsection{更新谱塑形}

E11 结果显示，Muon-like 极分解更新稳定地产生更平、更高秩的 update spectrum。在 noisy mini-batch default setting 的 equal-update control 下，nrUpdate 比值约为 \(2.017\)，stUpdate 比值约为 \(4.765\)。这说明极分解更新确实改变更新矩阵奇异值分布，但这只是机制事实，不是性能声明。

\begin{table}[t]
\centering
\caption{E11 已有局部几何证据摘要。比值均为 Muon-like/Adam 或 flat-polar/GD-spectrum。}
\label{tab:e11-core}
\begin{tabular}{llll}
\toprule
证据项 & 指标 & 数值 & 解释 \\
\midrule
更新谱塑形 & nrUpdate & $2.017\,[1.905,2.136]$ & Muon-like 更新核秩更高 \\
更新谱塑形 & stUpdate & $4.765\,[4.582,4.955]$ & Muon-like 更新稳定秩更高 \\
一阶校准 & $\rho(\Delta L,\langle G,D\rangle)$ & $0.9803\,[0.9787,0.9817]$ & one-step loss 由一阶内积很好解释 \\
Frobenius 预算 & flat/GD & $0.6071\,[0.5846,0.6304]$ & flat/polar 谱分配输给梯度谱分配 \\
Operator 预算 & flat/GD & $1.689\,[1.614,1.768]$ & flat/polar 谱分配胜出 \\
\bottomrule
\end{tabular}
\end{table}


\subsection{一阶损失下降校准}

E11 记录的 \texttt{update\_grad\_inner} 对应本文记号中的
\[
\langle G,D\rangle,
\]
即正的一阶下降 proxy。已有结果显示，one-step 实际损失下降与该 proxy 的 Spearman 相关约为 \(0.9803\)，说明在当前短步长范围内，使用局部一阶模型是合理的。

\subsection{范数几何边界}

谱分配 probe 固定梯度奇异向量，只改变奇异值分配，并分别在 Frobenius 预算与 operator-norm 预算下比较 flat/polar 与 GD-spectrum。结果表明：
\[
\text{Frobenius budget:}\quad \text{flat/GD}=0.6071<1,
\]
但
\[
\text{operator budget:}\quad \text{flat/GD}=1.689>1.
\]
这与附录~\ref{sec:appendix-proofs}中的一阶最优性推导完全一致。

\subsection{神经网络负控}

神经网络 sanity checks 显示，高 update rank 并不自动带来更好一阶进展。Deep MNIST MLP、MNIST Patch surrogate 与 MNIST ConvNet 都保留 Muon-like 更新的高核秩/高稳定秩 signature，但 one-step 对齐比仍显著低于 Adam。见表~\ref{tab:e11-neural-negative}。

\begin{table}[t]
\centering
\caption{神经网络 sanity checks：高秩更新不自动意味着更好一阶进展。}
\label{tab:e11-neural-negative}
\begin{tabular}{llll}
\toprule
任务 & nrUpdate & stUpdate & one-step 对齐比 \\
\midrule
Deep MNIST MLP & $2.262\,[2.175,2.353]$ & $18.31\,[16.26,20.62]$ & $0.5978\,[0.5633,0.6345]$ \\
MNIST Patch + 共享权重 & $4.235\,[3.834,4.677]$ & $12.89\,[11.77,14.12]$ & $0.5222\,[0.4943,0.5516]$ \\
MNIST ConvNet & $3.718\,[3.451,4.006]$ & $12.03\,[11.03,13.12]$ & $0.677\,[0.6321,0.7251]$ \\
\bottomrule
\end{tabular}
\end{table}


这部分证据支持本文的谨慎表述：
\[
\boxed{\text{SpecGrad 改变更新谱几何，但谱扩散是否有利取决于 activation/downstream geometry。}}
\]

\subsection{预测边界尚未完成}

当前边界预测器只能作为描述性分析。已有 audit 显示 in-sample fit 具有探索价值，但 held-out generalization 还不够稳定；部分 held-out setting 还存在 positive/negative label 退化问题。因此，当前不应声称已有一个可迁移的强预测规则。

\section{实验计划与待完成实验}
\label{sec:experiments-tbd}

本节区分已完成实验与待完成实验。当前完成的 E11 局部几何实验见第~\ref{sec:e11} 节。下面列出的实验尚未完成，本文暂时保留空位。

\subsection{实验一：Activation perturbation diagnostics}

\paragraph{状态.} \TBD。

\paragraph{目标.} 直接测量 SpecGrad / polar update 是否在 activation perturbation 上表现出 operator-norm 意义下的稳定性。

对每层记录
\[
\Delta Z_{\ell,\mathrm{direct}}
=(W_\ell^{t+1}-W_\ell^t)A_{\ell-1}^t.
\]
测量：
\[
\frac{\|\Delta Z_{\ell,\mathrm{direct}}\|_F}{\|Z_\ell^t\|_F},
\qquad
\frac{\|\Delta Z_{\ell,\mathrm{direct}}\|_{\op}}{\|Z_\ell^t\|_{\op}},
\]
以及 per-sample 最大扰动：
\[
\max_b
\frac{\|(W_\ell^{t+1}-W_\ell^t)h_{\ell-1}(x_b)\|_2}
{\|h_{\ell-1}(x_b)\|_2}.
\]

\paragraph{预期.} SpecGrad 应更自然地控制 operator/per-sample 指标，而不一定降低 batch-Frobenius movement。

\paragraph{结果.} \TBD。

\subsection{实验二：Activation stable rank 与层级条件}

\paragraph{状态.} \TBD。

计算
\[
\srank(A_{\ell-1})=\frac{\|A_{\ell-1}\|_F^2}{\|A_{\ell-1}\|_{\op}^2},
\qquad
\nrank(G_\ell)=\frac{\|G_\ell\|_*^2}{\|G_\ell\|_F^2}.
\]
检验条件
\[
\nrank(G_\ell)>\srank(A_{\ell-1})
\]
是否预测该层 polar/SpecGrad 的一阶优势。

\paragraph{结果.} \TBD。

\subsection{实验三：Downstream-aware activation perturbation}

\paragraph{状态.} \TBD。

目标是测量
\[
\|B_\ell D_\ell A_{\ell-1}\|_F
\]
或使用 JVP 估计
\[
\|J_\theta D_\ell\|_F.
\]
比较
\[
D_{\mathrm{spec}}=\operatorname{polar}(G_\ell)
\]
与
\[
D_{\mathrm{gd}}=G_\ell/\|G_\ell\|_F.
\]
核心指标为
\[
\frac{\langle G_\ell,D_{\mathrm{spec}}\rangle^2/\|J_\theta D_{\mathrm{spec}}\|_F^2}
{\langle G_\ell,D_{\mathrm{gd}}\rangle^2/\|J_\theta D_{\mathrm{gd}}\|_F^2}.
\]

\paragraph{结果.} \TBD。

\subsection{实验四：长尾小 batch 激活实验}

\paragraph{状态.} \TBD。

构造长尾分类或 synthetic long-tail mode 数据。检验 rare direction 可见性条件
\[
\frac{Bp_c}{1-\beta}\gtrsim 1.
\]
记录 tail-class loss、tail activation drift、tail activation coverage、minibatch spectral-tail consistency 与 effective nuclear rank。

\paragraph{结果.} \TBD。

\subsection{实验五：现代架构激活扰动 benchmark}

\paragraph{状态.} \TBD。

在 Transformer、现代 MLP 或 CNN 架构中测量
\[
\|\Delta A_\ell\|_F,
\qquad
\|\Delta A_\ell\|_{\op},
\qquad
\srank(A_\ell),
\qquad
\|J_\theta D_\ell\|_F.
\]
本实验不以 attention 为核心，而是直接观察 activation perturbation。

\paragraph{结果.} \TBD。

\section{讨论}

本文主张将谱梯度更新理解为 activation perturbation geometry，而不是简单理解为“更高 rank 的更新”。这一表述有三个好处。

第一，它解释了为什么 polar/SpecGrad 更新在 operator-norm 预算下自然。由于
\[
\|UV^\top\|_{\op}=1,
\]
它给出了清晰的 per-sample activation perturbation 控制：
\[
\|UV^\top a\|_2\le \|a\|_2.
\]

第二，它解释了为什么高秩更新可能失败。Batch-level 扰动由
\[
\|D A\|_F
\]
决定，而 polar 更新的上界是 \(\|A\|_F\)，可能随着 activation stable rank 增大而增大。因此，高秩谱扩散只有在这些方向带来足够一阶收益、并且 downstream sensitivity 不高时才有利。

第三，它自然引入 downstream-aware stable rank。若扰动通过后续网络传播为
\[
BDA,
\]
则相关几何量应为
\[
\ssrank(B,A),
\]
而非仅 \(\srank(A)\)。这说明 layerwise activation stable rank 可能是保守上界。

已有 E11 实验与这一观点一致：Muon-like 更新稳定地产生更平、更高秩的 update spectra，但进展优势是条件化的。当前证据不支持“Muon 一般更稳定”或“更高 rank 一般更好”的宽泛结论；更安全的说法是，Muon/SpecGrad 是 update-spectrum shaping 方法，其效应由任务、层、范数几何和下游敏感度共同决定。

\section{结论}

本文将谱梯度更新整理为一种激活扰动几何。核心结论如下：

\begin{enumerate}[label=(\arabic*)]
  \item Polar/SpecGrad 更新是 operator-norm 预算下的一阶最优方向，而不是 Frobenius 预算下的一阶最优方向。
  \item 对层激活而言，关键对象是 \(D A\)。SpecGrad 控制 operator/per-sample activation perturbation，但不保证 batch-Frobenius perturbation 更小。
  \item 对深层网络而言，真正相关的是 \(BDA\)。对应的敏感度比是
  \[
  \ssrank(B,A)=
  \frac{\sum_i \sigma_i(B)^2\sigma_i(A)^2}
  {\sigma_1(B)^2\sigma_1(A)^2}.
  \]
  \item 已有 E11 实验支持局部谱几何机制：Muon-like 更新重塑 update spectrum；one-step 下降由 \(\langle G,D\rangle\) 很好解释；polar 谱分配在 operator-norm 预算下有利，但在 Frobenius 预算下不利。
  \item 尚需补充直接的 activation perturbation、downstream Jacobian 和长尾小 batch 实验。
\end{enumerate}

因此，最稳妥的论文级表述是：
\[
\boxed{\text{谱梯度更新的收益不是来自高秩本身，而来自高秩谱扩散与激活/下游几何之间的匹配。}}
\]


\appendix
\section{核心证明细节}
\label{sec:appendix-proofs}

\subsection{谱范数球与核范数对偶}

设 \(G=U\Sigma V^\top\)。由 von Neumann trace inequality，任意 \(D\) 满足
\[
\langle G,D\rangle
\le
\sum_i\sigma_i(G)\sigma_i(D).
\]
若 \(\|D\|_{\op}\le \eta\)，则 \(\sigma_i(D)\le \eta\)，所以
\[
\langle G,D\rangle
\le
\eta\sum_i\sigma_i(G)
=\eta\|G\|_*.
\]
取 \(D=\eta UV^\top\) 达等号。因此
\[
\max_{\|D\|_{\op}\le\eta}\langle G,D\rangle=\eta\|G\|_*.
\]

\subsection{Frobenius 球与 Frobenius 范数对偶}

由 Cauchy--Schwarz，
\[
\langle G,D\rangle
\le
\|G\|_F\|D\|_F.
\]
若 \(\|D\|_F\le r\)，则
\[
\langle G,D\rangle\le r\|G\|_F.
\]
取 \(D=rG/\|G\|_F\) 达等号。

\subsection{Sandwiched sensitivity 的 Frobenius 情况}

对 \(J(D)=BDA\)，有
\[
\vecop(BDA)=(A^\top\otimes B)\vecop(D).
\]
因此
\[
\sup_{\|D\|_F\le1}\|BDA\|_F
=
\|A^\top\otimes B\|_{\op}.
\]
Kronecker product 的 operator norm 满足
\[
\|A^\top\otimes B\|_{\op}
=
\|A\|_{\op}\|B\|_{\op}.
\]
所以
\[
K_F^2=\|B\|_{\op}^2\|A\|_{\op}^2.
\]

\subsection{Sandwiched sensitivity 的 operator-norm 情况}

令
\[
B=U_B\Sigma_BV_B^\top,
\qquad
A=U_A\Sigma_AV_A^\top.
\]
则
\[
\|BDA\|_F=\|\Sigma_B Z\Sigma_A\|_F,
\qquad
Z=V_B^\top D U_A.
\]
若 \(\|D\|_{\op}\le1\)，则 \(\|Z\|_{\op}\le1\)。因此
\[
\sup_{\|D\|_{\op}\le1}\|BDA\|_F^2
=
\sup_{\|Z\|_{\op}\le1}\|\Sigma_BZ\Sigma_A\|_F^2.
\]
展开为
\[
\|\Sigma_BZ\Sigma_A\|_F^2
=
\Tr(\Sigma_B^2 Z\Sigma_A^2Z^\top).
\]
由重排不等式，最大值为
\[
\sum_i\sigma_i(B)^2\sigma_i(A)^2.
\]
取使 \(A\) 与 \(B\) 的奇异方向对齐的 partial isometry 可达等号。

\section{建议实验指标与记录协议}

后续实验建议记录以下量。

\subsection{更新谱指标}
\[
\nrank(D)=\frac{\|D\|_*^2}{\|D\|_F^2},
\qquad
\srank(D)=\frac{\|D\|_F^2}{\|D\|_{\op}^2}.
\]

\subsection{激活扰动指标}
\[
\Delta Z_{\ell,\mathrm{direct}}=(W_\ell^{t+1}-W_\ell^t)A_{\ell-1}^t.
\]
记录
\[
\|\Delta Z_{\ell,\mathrm{direct}}\|_F,
\qquad
\|\Delta Z_{\ell,\mathrm{direct}}\|_{\op},
\]
以及归一化版本。

\subsection{一阶进展指标}
\[
\langle G_\ell,D_\ell\rangle,
\qquad
\sum_\ell\langle G_\ell,D_\ell\rangle.
\]

\subsection{下游扰动指标}
用 JVP 估计
\[
\|J_\theta D\|_F^2.
\]
比较
\[
\frac{\langle G,D\rangle^2}{\|J_\theta D\|_F^2}
\]
在不同更新方向上的差异。

\section{安全表述与需避免表述}

\subsection{可以安全使用的表述}

\begin{itemize}
  \item SpecGrad / Muon-like 更新是 update-spectrum shaping 方法。
  \item Polar 更新是一阶线性化目标在 operator-norm 预算下的最优方向。
  \item 谱更新是否有利取决于任务、层、激活稳定秩、下游 Jacobian 敏感度和更新预算几何。
  \item 高 update rank 是机制事实，不是性能保证。
\end{itemize}

\subsection{应避免的表述}

\begin{itemize}
  \item Muon 一般比 Adam 更好。
  \item 更高 rank/stable rank 更新一般带来更低损失。
  \item 当前 E11 边界预测器已经能稳定预测 unseen tasks。
  \item 当前实验已经证明 SpecGrad 改善 activation stability 或 long-tail generalization。
\end{itemize}


\begin{thebibliography}{9}

\bibitem[Davis and Drusvyatskiy(2026)]{davis2026spectral}
Damek Davis and Dmitriy Drusvyatskiy.
\newblock \emph{When Do Spectral Gradient Updates Help in Deep Learning?}
\newblock arXiv:2512.04299, 2026.

\bibitem[Chen et~al.(2025)]{chen2025muon}
Lizhang Chen, Jonathan Li, and Qiang Liu.
\newblock \emph{Muon Optimizes Under Spectral Norm Constraints}.
\newblock arXiv:2506.15054, 2025.

\bibitem[Sanyal et~al.(2020)]{sanyal2020stable}
Amartya Sanyal, Philip H. S. Torr, and Puneet K. Dokania.
\newblock \emph{Stable Rank Normalization for Improved Generalization in Neural Networks and GANs}.
\newblock arXiv:1906.04659, 2020.

\bibitem[Liu(2026)]{e11repo}
Tianyang Liu.
\newblock \emph{E11 Condition Geometry Experiments for Muon Update-Spectrum Shaping}.
\newblock GitHub branch \texttt{codex/e11-three-problem-condition-geometry}, 2026.

\end{thebibliography}


\end{document}
